\documentclass[11pt]{article}
\usepackage[margin=28mm]{geometry}
\usepackage{amsmath,amssymb,amsthm,mathtools,mathrsfs,booktabs,longtable,array}
\usepackage[colorlinks]{hyperref}
\newcommand{\statement}[3]{\subsection*{#1. #3 [#2]}}
\begin{document}
\statement{C.31}{Proved}{Every full braid-operad endomorphism has a unit multiplier}
Let $\Phi$ be an object-fixing continuous endomorphism of the symmetric
operad $\widehat{\mathrm{PaB}}$ of full profinite parenthesized braids.
Let $\lambda\in1+2\widehat{\mathbb Z}$ be its exponent on the elementary
two-strand braiding. Then
\[
                         \lambda\in\widehat{\mathbb Z}^{\times}.
\]
Consequently the full Grothendieck--Teichm\"uller endomorphism monoid,
defined by the full duality, hexagon and pentagon equations with an odd
parameter, is already its group of units.

\begin{proof}
We first establish the precise full-completion input. Write $B_n$ and
$P_n$ for the ordinary and pure discrete Artin braid groups. The
finite-index inclusion $P_n\leq B_n$ induces its full profinite
topology, by taking cores of finite-index subgroups in $B_n$. Hence
\[
 1\longrightarrow\widehat P_n\longrightarrow\widehat B_n
       \longrightarrow S_n\longrightarrow1                     \tag{C.31.1}
\]
is exact with the full completion in the first term. Likewise the
split Fadell--Neuwirth sequence
$1\to F_{n-1}\to P_n\to P_{n-1}\to1$ retains the full topology on
its kernel. Indeed any finite-index subgroup of the finitely
generated free kernel contains a characteristic finite-index subgroup
$W$; the action on $F_{n-1}/W$ factors through a finite quotient of
$P_{n-1}$ and gives a finite semidirect-product quotient detecting
that kernel quotient.

Finite-rank free groups are good. Comparison of the discrete and
continuous Hochschild--Serre spectral sequences therefore gives
goodness of $P_n$ inductively, and then of $B_n$ using (C.31.1).
The finite-coefficient cohomology of the kernels in this induction is
finite, so the comparison applies to every required coefficient
module. These are also the precise primary-source conclusions of
\cite[Proposition 5.10 and Corollary 5.11]{CEHorel2017} and
\cite[\S2.3, Proposition 1]{CEMarin2012}.
The unordered configuration space of $n$ distinct points in
$\mathbb C$ is an aspherical smooth manifold of real dimension $2n$:
asphericity follows from the ordered Fadell--Neuwirth fibrations,
and passage to the free finite symmetric-group quotient preserves
the universal cover. Its CW dimension is at most $2n$, so
$\operatorname{cd}(B_n)\leq2n$. Goodness now implies, for every
finite continuous $p$-primary module $A$,
\[
 H^j_{\mathrm{cts}}(\widehat B_n,A)=0\qquad(j>2n).
\]
An arbitrary discrete $p$-primary module is the filtered union of
its finite invariant submodules, and continuous cochains on compact
profinite powers have finite image. Thus the same vanishing holds
for every such module and
$\operatorname{cd}_p(\widehat B_n)\leq2n$ for every prime $p$.
A closed $C_p$ is impossible, since cohomological dimension does
not increase on closed subgroups and
$H^j(C_p,\mathbb F_p)\ne0$ in every degree. In particular
\[
                     \widehat B_n\text{ is torsion-free}.        \tag{C.31.2}
\]
This implication is also recorded in
\cite[Lemma 8.4]{CEMorales2024}. It uses goodness and actual finite
cohomological dimension, rather than virtual cohomological dimension
or torsion-freeness of the discrete group alone.

Choose a binary parenthesization $b$ in arity $n$. The full subgroupoid
on its $S_n$-orbit, quotiented by the free label action of $S_n$, is
the one-object groupoid with automorphism group $B_n$. After full
completion its group is $\widehat B_n$, by (C.31.1). Since $\Phi$
fixes objects and commutes with this symmetric-group action, it
induces an actual continuous endomorphism
\[
                \phi_n:\widehat B_n\longrightarrow\widehat B_n
\]
whose map on $S_n$ is the identity. Taking the automorphism group of
just one labelled object would instead give the pure braid group;
the orbit quotient is essential here.

We next verify the exact full-twist formula without assuming that
$\phi_n$ is onto. Let $\vartheta$ be the one-strand ribbon twist.
The profinite ribbon-extension isomorphism
\cite[Proposition 8.1]{CEBHR2019} extends $\Phi$ uniquely to an
object-fixing endomorphism of $\widehat{\mathrm{PaRB}}$, preserving
the untwisted braids. The two-strand balancing calculation in
\cite[Proposition 7.7]{CEBHR2019}, valid over
$\widehat{\mathbb Z}$, gives
$\Phi(\vartheta)=\vartheta^\lambda$. More explicitly, if its
twist exponent is $t$, that calculation gives $2t=2\lambda$;
multiplication by $2$ is injective on $\widehat{\mathbb Z}$.
Put
\[
 z_n=(\sigma_1\cdots\sigma_{n-1})^n\in Z(B_n).
\]
The defining ribbon-cabling rule
\cite[Definition 6.4]{CEBHR2019} says literally
\[
 \vartheta\circ_1\operatorname{id}_b=(z_n,[1,\ldots,1]).
\]
Applying the ribbon endomorphism to the two sides gives respectively
$(z_n^\lambda,[\lambda,\ldots,\lambda])$ and
$(\phi_n(z_n),[\lambda,\ldots,\lambda])$. Consequently
\[
                         \phi_n(z_n)=z_n^\lambda.                \tag{C.31.3}
\]
This uses the operadic cabling identity; centrality of $z_n$ alone
would not justify (C.31.3) for an unproved automorphism.

Suppose now that $\lambda\in p\widehat{\mathbb Z}$ for some prime
$p$, and write $\lambda=p\nu$. In $B_p$ let
$\delta_p=\sigma_1\cdots\sigma_{p-1}$, so that
$\delta_p^p=z_p$. Centrality of $z_p$ persists in the full completion.
By (C.31.3) the element
\[
                  a=\phi_p(\delta_p)z_p^{-\nu}
\]
satisfies $a^p=1$. Its permutation is the nontrivial $p$-cycle of
$\delta_p$, since $z_p$ is pure and $\phi_p$ preserves permutations.
It therefore has exact order $p$, contradicting (C.31.2). Every
prime has been excluded, so $\lambda$ is a full unit.

Finally the full monoid is precisely the object-fixing operad
endomorphism monoid by \cite[Proposition 7.3]{CEBHR2019}. Its pairs
now have unit multiplier; C.30 makes their normalized free-group
endomorphisms automorphisms, and the closed-submonoid argument there
supplies inverses satisfying the same full equations.
\end{proof}

The use of arity $p$ for every prime is substantive. An action known
only in bounded arities would rule out only the corresponding
primes. Neither a pro-$p$ braid quotient nor a prounipotent action
supplies (C.31.1)--(C.31.2). In particular C.31 does not promote the
bare-pentagon element of A.37 to a full operad endomorphism: the full
hexagon and an actual profinite parameter still have to be supplied.

\end{document}
